% Figure: free response of a SDOF oscillator in the three damping
% regimes, released from x(0)=1 with zero initial velocity.
% Underdamped zeta=0.15, critically damped zeta=1, overdamped zeta=2.5,
% all on the same natural frequency omega_n = 1.
\begin{figure}[htbp]
\centering
\begin{tikzpicture}
\begin{axis}[
  width=12cm, height=7cm,
  xlabel={$\omega_n t$}, ylabel={$x/x_0$},
  xmin=0, xmax=18, ymin=-0.55, ymax=1.08,
  axis lines=left,
  domain=0:18, samples=400,
  legend pos=north east,
  legend cell align=left,
  every axis x label/.style={at={(current axis.right of origin)}, anchor=west},
  every axis y label/.style={at={(current axis.above origin)}, anchor=south},
  grid=major, grid style={enggray!20},
  tick label style={font=\scriptsize},
  label style={font=\small},
  legend style={font=\scriptsize},
]
% Underdamped, zeta = 0.15, wd = sqrt(1-0.0225) = 0.98869.
% x(0)=1, xdot(0)=0 -> x = e^{-z t}(cos wd t + (z/wd) sin wd t).
\addplot[engblue, very thick] {exp(-0.15*x)*(cos(deg(0.98869*x)) + (0.15/0.98869)*sin(deg(0.98869*x)))};
\addlegendentry{underdamped $\zeta=0.15$}
% Critically damped, zeta = 1: x = (1 + t) e^{-t}.
\addplot[englime, very thick] {(1 + x)*exp(-x)};
\addlegendentry{critical $\zeta=1$}
% Overdamped, zeta = 2.5: roots s = -z +- sqrt(z^2-1).
% s1 = -2.5 + sqrt(5.25) = -0.20871, s2 = -2.5 - 2.29129 = -4.79129.
% x(0)=1, xdot(0)=0: C1 = -s2/(s1-s2), C2 = s1/(s1-s2).
% s1-s2 = 4.58258, C1 = 4.79129/4.58258 = 1.04555, C2 = -0.04555.
\addplot[engred, very thick, dashed] {1.04555*exp(-0.20871*x) - 0.04555*exp(-4.79129*x)};
\addlegendentry{overdamped $\zeta=2.5$}
% zero line
\addplot[enggray, thin, forget plot] {0};
\end{axis}
\end{tikzpicture}
\caption[Free response in the three damping regimes]{Free response of
$m\ddot{x}+c\dot{x}+kx=0$ released from $x(0)=x_0$ with zero initial
velocity, for the same natural frequency $\omega_n$ and three damping
ratios. The underdamped case (blue, $\zeta=0.15$) overshoots through
zero and rings down inside an exponential envelope. The critically
damped case (green, $\zeta=1$) returns to equilibrium without crossing
zero and does so faster than any heavier damping. The overdamped case
(red dashed, $\zeta=2.5$) returns without overshoot but more slowly,
its long-time decay set by the slow root $s_1=-\zeta\omega_n+\omega_n
\sqrt{\zeta^2-1}$, which approaches zero as the damping grows. Critical
damping is the boundary that makes the characteristic discriminant
$\zeta^2-1$ vanish, the fastest non-overshooting return. Computed from
the closed-form solutions of each regime.}
\label{fig:vol03ch06:damping-regimes}
\end{figure}
